\diarytitle{0102}{arctan x の逆関数の微分法による導関数の導出}

$x = \tan y$（$-\dfrac{\pi}{2} < y < \dfrac{\pi}{2}$）とおいて両辺を $x$ で微分すると
\[
1 = \sec^{2} y \cdot \frac{dy}{dx} = (1+\tan^{2}y)\frac{dy}{dx} = (1+x^{2})\frac{dy}{dx}
\]
したがって
\[
\boxed{\; \frac{dy}{dx} = \frac{1}{1+x^{2}} \;}
\]
（検算: $x=1$ のとき $y = \arctan 1 = \dfrac{\pi}{4}$ で傾きは $\dfrac{1}{1+1}=\dfrac{1}{2}$。$\tan y$ の逆関数として、$x=1$ 付近で $\tan y$ の傾き $\sec^{2}(\pi/4)=2$ の逆数 $\dfrac12$ と一致する。）
